\section{Hypotheses}\label{sec:hypotheses}

The study tests four hypotheses. Each is examined through several measures, which are
reported as parts of the hypothesis rather than as separate tests.

\paragraph{H1: Expiry concentrates trading in the settlement window.} The closing and rolling
of expiring positions, and the unwinding of arbitrage positions at the settlement price, are
expected to draw volume into the window and to make it more expensive to trade in. The
measures are the window's share of session volume, the quoted spread, displayed depth, the
cost of sweeping the displayed book and the rate of cancellation
(Sections~\ref{sec:activity} and~\ref{sec:book}).

\paragraph{H2: Trading in the window is one-sided on expiry days.} Pressure on the settlement
price requires order flow that leans persistently in one direction. It is measured by the order
imbalance
\begin{equation*}
  \mathrm{OIB} = \frac{BI - SI}{BI + SI} \in [-1, 1],
\end{equation*}
where $BI$ and $SI$ are buyer- and seller-initiated volume, over the window and in each of its
five-minute blocks, by the persistence of the imbalance from one block to the next, and by the
price move from the start of the window to its end. A market with balanced two-sided trading
has an imbalance near zero that does not persist; one-sided pressure produces an imbalance
that stays away from zero (Section~\ref{sec:pressure}).

\paragraph{H3: The settlement price is displaced by trading in the window.} A settlement price
moved by trading that does not reflect information is reversed once that trading stops, and a
displacement caused by the settlement should appear only where something settles. The
measures are the reversal of the window's move the next morning, and the signed move of the
index, whose derivatives are settled in cash, relative to an index without derivatives
(Section~\ref{sec:distortion}).

\paragraph{H4: The derivatives create an incentive to move the settlement price.} The incentive
depends on how the derivative is settled. It is examined through the pricing of the expiring
stock future within the window and through the attraction of the settlement price to option
strikes, for stocks and for the index (Section~\ref{sec:derivatives}).

Each hypothesis is tested on monthly expiries against other sessions, with the weekly index
expiries and month ends as calendar controls, and in securities with derivatives against the
placebo group.
